\documentclass[UTF8,fontset=none,aspectratio=169,10pt]{beamer}
\usepackage{ads-beamer-theme}
\title{高级分布式系统}
\subtitle{第十三讲：Failure Detector 与不完美信息}
\author{陈鹏飞\quad 计算机学院}
\institute{中山大学}
\date{2025--2026 学年}

\begin{document}
\begin{frame}[plain]
\begin{tikzpicture}[remember picture,overlay]
  \fill[ADSBlue] (current page.north west) rectangle ([xshift=2.4cm]current page.south west);
  \fill[ADSOrange] ([xshift=2.4cm]current page.north west) rectangle ([xshift=2.55cm]current page.south west);
  \node[anchor=west,align=left] at ([xshift=3.3cm,yshift=1.0cm]current page.west) {\fontsize{32}{38}\selectfont\color{ADSBlue}\bfseries 高级分布式系统};
  \node[anchor=west,align=left] at ([xshift=3.3cm,yshift=-0.4cm]current page.west) {\fontsize{21}{27}\selectfont\color{ADSDarkBlue}第十三讲\quad Failure Detector};
  \node[anchor=west,align=left] at ([xshift=3.3cm,yshift=-2.2cm]current page.west) {\large\color{ADSGray}当“慢”与“坏掉”无法被准确区分};
  \node[anchor=east] at ([xshift=-0.8cm,yshift=-3.2cm]current page.east) {\small\color{ADSGray}陈鹏飞\quad 中山大学\quad 2025--2026};
  \draw[ADSBlue,line width=0.8pt] ([xshift=3.3cm,yshift=-1.35cm]current page.west) -- ([xshift=12.6cm,yshift=-1.35cm]current page.west);
\end{tikzpicture}
\end{frame}

\sectionframe{一、故障检测器不是水晶球}{它是一个分布式 oracle：提供“怀疑谁失联”的提示，但提示可能暂时错误。}

\begin{frame}{本讲学习目标}
\begin{enumerate}
  \item 解释同步与异步系统中“检测故障”的根本差异；
  \item 区分 completeness、accuracy 及其强/弱、最终版本；
  \item 理解 Failure Detector 类别和 reducibility；
  \item 说明故障检测器如何支撑共识、原子广播和服务切换；
  \item 把误判、超时和恢复映射到 Kubernetes 与 AI 集群。
\end{enumerate}
\vfill
\callout{主线}{检测器的输出不是事实本身，而是带安全/活性取舍的证据。}
\end{frame}

\begin{frame}{同步与异步系统}
\begin{tabularx}{\linewidth}{>{\bfseries}p{2.8cm}XX}
\toprule
同步系统 & 消息传输、进程执行和时钟误差都有已知上界；超时具有较强语义。 \\
异步系统 & 没有已知传输和处理上界；超时只能说明“暂时没看到”，不能证明崩溃。 \\
\bottomrule
\end{tabularx}
\vfill
\callout{关键难题}{在纯异步系统里，慢进程和崩溃进程对观察者不可区分。}
\end{frame}

\begin{frame}{两个不可同时完美满足的场景}
\begin{columns}[T]
\column{.47\linewidth}
\begin{block}{q 只是很慢}
P 误判 q 已崩溃并切换主节点，q 其实仍会继续写入，产生双主或重复提交。
\end{block}
\column{.47\linewidth}
\begin{block}{q 确实崩溃}
P 为避免误判而无限等待 q，整个系统的恢复与进度都被拖住。
\end{block}
\end{columns}
\vfill
\callout{结论}{安全性要求不要过早排除正确进程；活性要求最终要敢于怀疑。}
\end{frame}

\sectionframe{二、Completeness 与 Accuracy}{故障检测器的质量由“最终能否发现坏进程”和“是否误伤好进程”两条轴描述。}

\begin{frame}{Completeness：最终发现故障}
\begin{itemize}
  \item Strong Completeness：每个崩溃进程最终被每个正确进程永久怀疑；
  \item Weak Completeness：每个崩溃进程最终至少被某个正确进程永久怀疑；
  \item 强版本更便于全局协调，但需要更多传播；
  \item 弱版本可通过广播怀疑集合，模拟为强版本。
\end{itemize}
\end{frame}

\begin{frame}{Accuracy：尽量不怀疑正确进程}
\begin{itemize}
  \item Strong Accuracy：正确进程在真正崩溃前从不被怀疑；
  \item Weak Accuracy：至少有一个正确进程永远不被怀疑；
  \item Eventual Strong Accuracy：经过某个未知时刻后，正确进程不再被误判；
  \item Eventual Weak Accuracy：经过某个未知时刻后，至少一个正确进程不再被误判。
\end{itemize}
\vfill
\callout{工程现实}{网络抖动、GC、GPU 预热和模型加载都会让“最终准确”比“永久准确”更现实。}
\end{frame}

\begin{frame}{Failure Detector 类别}
\begin{center}
\begin{tikzpicture}[node distance=8mm and 11mm,>=Latex]
  \nodebox[minimum width=2.6cm]{ADSBlue}{pfd}{Perfect $P$\\强完整 + 强准确};
  \nodebox[right=of pfd,minimum width=2.6cm]{ADSOrange}{spfd}{Strong $S$\\强完整 + 弱准确};
  \nodebox[below=of pfd,minimum width=2.6cm]{ADSGreen}{wfd}{Eventually Perfect $\Diamond P$\\强完整 + 最终强准确};
  \nodebox[right=of wfd,minimum width=2.6cm]{ADSDarkBlue}{ewfd}{Eventually Strong $\Diamond S$\\强完整 + 最终弱准确};
\end{tikzpicture}
\end{center}
\begin{itemize}
  \item 类别不是实现名称，而是算法可依赖的性质集合；
  \item 异步 Crash 共识常以 $\Diamond S$ 作为足够的故障检测能力；这里的“最终”只修饰准确性，不修饰强完整性。
\end{itemize}
\end{frame}

\begin{frame}{心跳检测器的基本实现}
\begin{enumerate}
  \item 每个节点周期性发送 heartbeat 或 lease renewal；
  \item 接收者记录最近一次成功时间；
  \item 超过自适应阈值后输出 suspect；
  \item 之后若收到迟到 heartbeat，可以撤销怀疑或等待更高 epoch；
  \item 误判率、检测延迟和网络开销形成三方权衡。
\end{enumerate}
\end{frame}

\begin{frame}{为什么固定超时在云环境里危险？}
\begin{itemize}
  \item 节点启动、GPU 驱动加载、JVM GC、模型拉取都有长尾；
  \item 固定 timeout 把正常长尾误判为故障；
  \item 自适应 timeout 可基于 RTT 分位数、历史抖动和工作负载阶段调整；
  \item 但过度自适应会让真正故障被拖延，需要上限、租约和健康探针组合。
\end{itemize}
\end{frame}

\sectionframe{三、Reducibility：弱检测器如何支撑强算法}{如果能用一个检测器模拟另一个检测器，就可以复用已有共识和原子广播算法。}

\begin{frame}{Failure Detector 的可归约性}
\begin{itemize}
  \item 若存在算法 $T_{D\rightarrow D'}$，能用 D 的输出模拟 D'；
  \item 记作 $D\succeq D'$，表示 D 至少和 D' 一样强；
  \item 若双方互相可归约，则两者等价；
  \item 算法设计时可以先证明“需要哪个性质”，再选择实现可获得的检测器。
\end{itemize}
\end{frame}

\begin{frame}{弱完整性到强完整性}
\begin{enumerate}
  \item 每个进程查询自己的本地检测器，得到 suspects；
  \item 周期性把 suspects 集合广播给所有其他进程；
  \item 收到其他进程的 suspects 后取并集，排除消息发送者自己；
  \item 只要某个正确进程最终怀疑了崩溃者，所有正确进程都会得到该证据。
\end{enumerate}
\vfill
\callout{代价}{传播怀疑集合需要额外消息，并且集合可能膨胀；工程中常用 epoch、压缩位图或增量更新。}
\end{frame}

\begin{frame}{故障检测器与 leader 选举}
\begin{itemize}
  \item 检测器提供“哪些候选者暂时不可用”的提示；
  \item 选举协议使用优先级、epoch 和多数派确认选择 leader；
  \item 不能仅凭一个本地 suspect 就宣布新 leader，否则会产生双主；
  \item leader 的有效性通常需要租约、fencing token 或共识日志确认。
\end{itemize}
\end{frame}

\begin{frame}{误判恢复：suspect 不等于 delete}
\begin{itemize}
  \item Suspect：停止向节点发送新任务、降低流量、等待更多证据；
  \item Fenced：通过 epoch/token 使旧节点即使恢复也不能继续写入；
  \item Delete：从控制面移除节点对象并重新调度工作；
  \item 三者的不可逆程度不同，应该按证据逐步升级。
\end{itemize}
\end{frame}

\sectionframe{四、共识为何需要 Failure Detector}{纯异步系统中，即使只有一个 crash failure，也无法同时保证终止、完整性和一致性。}

\begin{frame}{共识的四个性质}
\begin{tabularx}{\linewidth}{>{\bfseries}p{2.9cm}X}
\toprule
Termination & 每个正确进程最终决定某个值。 \\
Uniform Integrity & 每个进程最多决定一次。 \\
Agreement & 不同正确进程不能决定不同值。 \\
Uniform Validity & 决定的值必须由某个进程提出。 \\
\bottomrule
\end{tabularx}
\vfill
\callout{联系前面}{互斥只选一个进入者；共识要让所有正确进程对同一个值形成不可分叉决定。}
\end{frame}

\begin{frame}{FLP 直觉：等待还是决定？}
\begin{itemize}
  \item 在异步系统中，下一条消息可能只是延迟，也可能永远不会到；
  \item 如果进程过早决定，可能破坏 Agreement；
  \item 如果进程永远等待，可能破坏 Termination；
  \item Failure Detector 提供不完美但最终有用的时间/故障信息，打破这种对称性。
\end{itemize}
\end{frame}

\begin{frame}{轮换协调者算法的直觉}
\begin{enumerate}
  \item 进程按轮次运行，若轮次从 1 开始，则协调者为 $((r-1)\bmod n)+1$；
  \item 协调者广播自己的估计值；
  \item 其他进程报告已知估计和 suspect 信息；
  \item 若协调者未崩溃且最终不再被怀疑，估计被锁定并决定；
  \item 若协调者失败，进入下一轮。
\end{enumerate}
\vfill
\callout{关键机制}{最终准确性保证“迟早有一轮协调者不会再被正确进程怀疑”。}
\end{frame}

\begin{frame}{轮换协调者完整例子：第 1 轮协调者崩溃}
\small
三个进程的初始估计为 $P_1:(A,0)$、$P_2:(B,0)$、$P_3:(B,0)$；第二项是估计值最近被协调者采纳的轮次。
\begin{center}
\begin{tabularx}{.96\linewidth}{>{\bfseries}p{1.4cm}>{\raggedright\arraybackslash}p{3.2cm}>{\raggedright\arraybackslash}X}
\toprule
步骤 & 消息/故障 & 状态变化 \\
\midrule
1 & 第 1 轮协调者为 $P_1$ & $P_1$ 提议 $PROPOSE(1,A)$。 \\
2 & 提议只到达 $P_2$ & $P_2$ 更新估计为 $(A,1)$ 并回复 ACK。 \\
3 & $P_1$ 随后崩溃 & 未收集多数派 ACK，也没有产生 DECIDE；$P_3$ 仍保持 $(B,0)$。 \\
4 & 检测器怀疑 $P_1$ & $P_2,P_3$ 不把怀疑结果当作决定，而是进入第 2 轮。 \\
\bottomrule
\end{tabularx}
\end{center}
\callout{关键点}{“收到协调者提议”不等于“已经决定”；没有多数派确认，第 1 轮只能留下新的估计，不能留下决定。}
\end{frame}

\begin{frame}{轮换协调者完整例子：第 2 轮完成决定}
\small
第 2 轮协调者为 $P_2$，多数派大小为 2。
\begin{center}
\begin{tabularx}{.96\linewidth}{>{\bfseries}p{1.3cm}>{\raggedright\arraybackslash}p{3.4cm}>{\raggedright\arraybackslash}X}
\toprule
步骤 & 消息 & 作用 \\
\midrule
1 & $P_2$ 收集 STATE & 从 $P_2$ 得到 $(A,1)$，从 $P_3$ 得到 $(B,0)$。 \\
2 & 选择最高轮次估计 & $P_2$ 选择 $A$，广播 $PROPOSE(2,A)$。 \\
3 & $P_2,P_3$ 回复 ACK & 第 2 轮形成多数派确认；旧的第 1 轮消息不能覆盖新轮次状态。 \\
4 & $P_2$ 广播 $DECIDE(A)$ & 所有正确进程最终决定 $A$。 \\
\bottomrule
\end{tabularx}
\end{center}
\begin{itemize}
  \item 若更早轮次已经形成多数派锁定，新多数派必与它相交，并携带最高轮次的锁定值；
  \item 最终准确性保证迟早有一轮的正确协调者不再被误判，从而完成决定。
\end{itemize}
\end{frame}

\sectionframe{五、AI 集群中的故障检测}{心跳只是起点；AI 任务需要区分进程、GPU、网络、存储和训练进度故障。}

\begin{frame}{多层健康信号}
\begin{tabularx}{\linewidth}{>{\bfseries}p{2.8cm}XX}
\toprule
进程层 & liveness/readiness probe、心跳、线程堆栈 & 进程是否运行 \\
设备层 & GPU ECC、显存、温度、XID、NVLink 错误 & 加速器是否可靠 \\
网络层 & RTT、丢包、RDMA counter、collective timeout & 通信是否可完成 \\
作业层 & step time、loss、checkpoint age、队列进度 & 训练是否真正前进 \\
\bottomrule
\end{tabularx}
\vfill
\callout{解释}{一个 Pod 处于 Running，不代表训练作业没有逻辑死锁或梯度通信卡住。}
\end{frame}

\begin{frame}{训练作业中的误判例子}
\begin{enumerate}
  \item GPU 节点正在拉取 20GB 模型权重，30 秒内没有应用心跳；
  \item 控制器把节点判为 failed，重新调度同一作业；
  \item 两份权重同时加载，存储和网络压力更大；
  \item 正确策略是分阶段 readiness：节点可达、运行时就绪、模型加载、warm-up 完成。
\end{enumerate}
\end{frame}

\begin{frame}{课堂互动：suspect 后能不能立即抢占？}
\begin{itemize}
  \item 场景：训练 worker 90 秒无心跳，但 GPU 利用率 95\%，网络仍有流量；
  \item 问题 1：这是进程故障、探针故障还是长尾阶段？
  \item 问题 2：应该先降级流量、发送 fencing token，还是立即杀 Pod？
  \item 问题 3：如何避免旧 worker 恢复后继续写 checkpoint？
\end{itemize}
\vfill
\callout{推荐回答}{先收集多层证据，再用 epoch/fencing 隔离旧实例，最后决定是否抢占。}
\end{frame}

\sectionframe{六、故障检测器的完整决策链}{从心跳缺失到安全降级，中间至少有 suspect、fence 和恢复确认。}

\begin{frame}{一个 3 节点时间线}
\begin{center}
\begin{tikzpicture}[>=Latex,x=0.15cm]
\draw[ADSGray,->] (0,0)--(72,0) node[right,font=\scriptsize]{现实时间};
\foreach \t/\lab in {0/t0,10/心跳,20/心跳,30/延迟,40/超时,52/恢复,64/续租}{\draw[ADSGray] (\t,.1)--(\t,-.1) node[below,font=\tiny]{\lab};}
\draw[ADSBlue,very thick] (10,.35)--(10,-.35) node[above,font=\scriptsize]{B 正常};
\draw[ADSOrange,very thick] (40,.35)--(40,-.35) node[above,font=\scriptsize]{suspect};
\draw[ADSOrange,dashed,very thick] (45,.35)--(45,-.35) node[above,font=\scriptsize]{fence};
\draw[ADSGreen,very thick] (64,.35)--(64,-.35) node[above,font=\scriptsize]{新 epoch};
\end{tikzpicture}
\end{center}
\begin{itemize}
  \item B 在 30--40 之间只是“没有及时响应”，不等于已经崩溃；
  \item suspect 触发保护动作，fencing 阻止旧实例继续写；
  \item B 恢复后必须用新 epoch 重新加入，而不是直接接管旧身份。
\end{itemize}
\end{frame}

\begin{frame}{Completeness 与 Accuracy 的四种组合}
\begin{tabularx}{\linewidth}{>{\bfseries}p{3.2cm}XX}
\toprule
强完整 & 所有崩溃进程最终被怀疑。 & 不会永久漏报，但可能误报。 \\
弱完整 & 至少一个正确进程最终怀疑每个崩溃进程。 & 常用于选主。 \\
强准确 & 永不怀疑正确进程。 & 异步网络中通常不可能免费得到。 \\
最终强准确 & 错误怀疑会停止，最终只保留真实故障。 & 实际系统常用的折中。 \\
\bottomrule
\end{tabularx}
\end{frame}

\begin{frame}{超时阈值的数值例子}
假设过去 1000 个心跳 RTT 的 P99 为 180ms：
\begin{itemize}
  \item timeout=200ms：检测快，但网络抖动时误报多；
  \item timeout=1s：误报少，但故障期间多等待 800ms；
  \item 自适应 timeout：$\max(\text{P99}\times k,\text{最小安全窗口})$；
  \item 训练集群和在线服务应使用不同的阈值与降级策略。
\end{itemize}
\vfill
\callout{实验}{把延迟分布从均值改成长尾分布，比较固定阈值和自适应阈值的误判率。}
\end{frame}

\begin{frame}{为什么 lease 不能替代 fencing？}
\begin{enumerate}
  \item 旧 leader 的本地时钟认为 lease 尚未过期；
  \item 新 leader 已经获得控制权；
  \item 如果存储只检查时间，两个 leader 都可能写入；
  \item fencing token 用单调 epoch 拒绝旧 token，即使旧 leader 恢复也不能写。
\end{enumerate}
\end{frame}

\begin{frame}{Kubernetes：Running、Ready、Lease}
\begin{center}
\begin{tikzpicture}[>=Latex,node distance=12mm]
\nodebox[minimum width=2.2cm]{ADSBlue}{run}{Running};
\nodebox[right=of run,minimum width=2.2cm]{ADSOrange}{ready}{Ready};
\nodebox[right=of ready,minimum width=2.2cm]{ADSGreen}{lease}{NodeLease};
\nodebox[below=of ready,minimum width=2.5cm]{ADSGray}{fence}{驱逐/隔离};
\draw[\arrow,ADSGray] (run)--node[above,font=\scriptsize]{进程存在}(ready);
\draw[\arrow,ADSGray] (ready)--node[above,font=\scriptsize]{节点续租}(lease);
\draw[\arrow,ADSOrange] (lease)--(fence);
\end{tikzpicture}
\end{center}
\begin{itemize}
  \item Running 只说明容器进程存在；
  \item Ready 说明当前可以接收工作；
  \item Lease 说明节点仍与控制面保持联系；
  \item 三个谓词失败后的动作不同，不能用一个 timeout 统一替代。
\end{itemize}
\end{frame}

\begin{frame}{AI 训练：检测“活着但不推进”}
\begin{itemize}
  \item worker 仍能发送心跳，但 step time 从 2s 升到 30s；
  \item GPU 利用率下降、NCCL 等待上升，说明逻辑活性变差；
  \item detector 应同时观察心跳、step progress、checkpoint age 和队列变化；
  \item 只有进程故障才重启，逻辑不推进可以先采样 trace、降级或迁移。
\end{itemize}
\end{frame}

\begin{frame}{故障检测器的审计日志}
\begin{tabularx}{\linewidth}{>{\bfseries}p{3cm}XX}
\toprule
字段 & 例子 & 用途 \\
subject & node-B / leader-3 & 明确检测对象 \\
evidence & heartbeat timeout / P99 anomaly & 解释为何怀疑 \\
epoch & 42 & 防止旧实例写入 \\
action & suspect / fence / restart & 区分可逆和破坏性动作 \\
recovery & lease renewed / state rebuilt & 记录恢复闭环 \\
\bottomrule
\end{tabularx}
\end{frame}

\begin{frame}{课堂推演：一次误判如何传播？}
\begin{enumerate}
  \item B 只是 GC 停顿 3 秒，A 将其标为 suspect；
  \item 控制器把 B 从服务后端摘除，并启动替代实例 C；
  \item B 恢复后仍持有旧缓存和旧写 token；
  \item 讨论：哪些动作需要 fencing？哪些动作可以撤销？如何避免双写？
\end{enumerate}
\vfill
\callout{答案关键词}{suspect 不等于 delete；旧 epoch 必须失效；缓存恢复需要版本校验。}
\end{frame}

\begin{frame}{本讲总结}
\begin{enumerate}
  \item 异步系统中无法准确区分“慢”和“坏”，Failure Detector 只能提供提示；
  \item Completeness 关注发现崩溃，Accuracy 关注不误伤正确进程；
  \item 强/弱、最终强/弱性质组成检测器类别，并可通过 reducibility 复用算法；
  \item 共识、leader 选举和原子广播依赖最终有用的故障信息；
  \item AI 集群必须组合进程、GPU、网络、存储和作业进度信号。
\end{enumerate}
\vfill
\callout{下一讲衔接}{下一讲把故障检测器放进完整的 Consensus、Raft/Paxos 和日志复制流程。}
\end{frame}

\begin{frame}{超时阈值如何设置？}
\begin{itemize}
  \item 低阈值：检测快，但误判多；
  \item 高阈值：误判少，但故障恢复慢；
  \item 可按 RTT 的 P99/P999、GC 暂停、模型加载阶段动态调整；
  \item 阈值更新要有上下限和冷却时间，避免探测器自身振荡。
\end{itemize}
\vfill
\callout{课堂计算}{给出 RTT 样本和故障恢复预算，让学生选择 timeout、heartbeat 周期与重试次数。}
\end{frame}

\begin{frame}{检测器输出的生命周期}
\begin{center}
\begin{tikzpicture}[node distance=9mm and 12mm,>=Latex]
  \nodebox{ADSBlue}{alive}{Alive};
  \nodebox[right=of alive]{ADSOrange}{suspect}{Suspect};
  \nodebox[right=of suspect]{ADSGreen}{fenced}{Fenced};
  \nodebox[right=of fenced]{ADSDarkBlue}{dead}{Removed};
  \draw[\arrow,ADSGray] (alive)--(suspect);
  \draw[\arrow,ADSGray] (suspect)--(fenced);
  \draw[\arrow,ADSGray] (fenced)--(dead);
  \draw[->,ADSGray] (suspect.south) .. controls +(0,-1) and +(0,-1) .. node[below,font=\scriptsize]{迟到心跳}(alive.south);
\end{tikzpicture}
\end{center}
\begin{itemize}
  \item 不是所有 suspect 都需要升级为 removed；
  \item fencing 让恢复的旧实例失去写权限，是避免双主的关键。
\end{itemize}
\end{frame}

\begin{frame}{Leader 租约与故障检测}
\begin{itemize}
  \item leader 周期性续租 lease；
  \item lease 到期后，其他候选者可以尝试获取新 epoch；
  \item 旧 leader 即使网络恢复，也必须被 fencing token 拒绝写入；
  \item 时间租约需要时钟假设，epoch/fencing 更适合处理时钟漂移和网络分区。
\end{itemize}
\end{frame}

\begin{frame}{故障检测与原子广播}
\begin{enumerate}
  \item 原子广播要保证所有正确进程以相同顺序交付；
  \item leader 负责排序和广播，故障检测器帮助识别 leader 是否失效；
  \item suspicion 不能直接改变交付顺序，必须通过新 epoch 或共识确认；
  \item 因而“怀疑”是触发协议转换的输入，而不是最终决定。
\end{enumerate}
\end{frame}

\begin{frame}{Kubernetes 探针与 Failure Detector 的对应}
\begin{tabularx}{\linewidth}{>{\bfseries}p{3cm}XX}
\toprule
livenessProbe & 进程是否需要重启 & 类似“本地进程是否存活” \\
readinessProbe & 是否加入服务后端 & 类似“当前是否可提供工作” \\
startupProbe & 是否仍在启动阶段 & 避免长启动被误杀 \\
NodeLease & 节点是否持续续租 & 提供节点级心跳/租约 \\
\bottomrule
\end{tabularx}
\vfill
\callout{解释}{Running、Ready 和可被调度不是同一个谓词，不能用一个探针替代全部语义。}
\end{frame}

\begin{frame}{AI 训练的“活性”指标}
\begin{itemize}
  \item step time 持续增长：可能是通信退化或数据 I/O 瓶颈；
  \item loss 长时间不变：可能是训练数值问题，也可能是 worker 卡住；
  \item checkpoint age 过大：说明作业虽然存活，但恢复窗口正在变差；
  \item 队列中的请求不减少：说明服务 Ready 但逻辑工作没有推进。
\end{itemize}
\end{frame}

\begin{frame}{课堂小实验：构造一次误判}
\begin{enumerate}
  \item 模拟一个 3 节点集群，设置消息延迟在 50ms--5s 之间变化；
  \item 让节点 B 在 2s 内暂停处理，但没有崩溃；
  \item 比较固定 timeout、P99 timeout 和 lease + fencing 三种策略；
  \item 记录误判次数、恢复延迟、重复写入和控制消息量。
\end{enumerate}
\end{frame}

\begin{frame}{设计检查清单}
\begin{enumerate}
  \item 检测器提供的是哪种 completeness/accuracy？
  \item 误判时，系统会丢数据、双主，还是只暂时降级？
  \item suspicion 到 fencing/remove 之间是否有可逆阶段？
  \item 恢复的旧实例如何被拒绝写入？
  \item 故障检测器本身的状态如何持久化和扩展？
\end{enumerate}
\end{frame}

\begin{frame}{课后题}
\begin{enumerate}
  \item 分别给出强完整、弱完整、强准确和最终强准确的例子；
  \item 设计一个 Kubernetes 节点多层 Failure Detector；
  \item 说明为什么 suspect 不能直接等价为 delete；
  \item 用轮换协调者解释故障检测器如何帮助共识终止。
\end{enumerate}
\end{frame}

\begin{frame}{故障检测器的边界：怀疑不是决定}
\small
\begin{itemize}
  \item 在异步系统中，“超时未收到心跳”同时可能表示崩溃、网络拥塞、GC 暂停或调度延迟；suspect 只是一个可撤销的估计。
  \item Kubernetes 的 Running、Ready、Lease 和 NodeCondition 是不同层次的信号，不能拼成一个天然线性的“正常—故障—删除”状态机。
  \item 只有共识、epoch/fencing 或持久化任期等机制，才能把估计转换成唯一的领导权和写入权；收到 coordinator proposal 也不等于已经决定。
  \item 设计恢复动作时，应先隔离旧实例的副作用，再清理资源；否则误判可能造成双主、重复写入或两个 controller 同时调度。
\end{itemize}
\vfill
\callout{安全优先}{故障检测器可以帮助系统推进，但不能单独创造“某节点已经永久失败”的事实。}
\end{frame}

\end{document}
